\documentclass[sigconf]{acmart}
\setcopyright{none}
\settopmatter{printacmref=false,printfolios=false}
\renewcommand\footnotetextcopyrightpermission[1]{}
\acmDOI{}
\acmISBN{}
\acmConference[SENG 474]{SENG 474 Evaluation}{University of Victoria}{Canada}
\usepackage{booktabs}
\usepackage{graphicx}
\begin{document}
\title{Evaluating a Convolutional Surrogate for Lid-Driven Cavity Flow}
@@AUTHORS@@
\begin{abstract}
We evaluate a convolutional neural network that predicts the next horizontal
velocity, vertical velocity, and pressure fields of a numerical cavity solver.
We compare it with persistence and a local ridge regressor using independent
simulation trajectories, grouped cross-validation, and a reserved final test set.
One-step errors and autoregressive rollouts measure predictive accuracy and
error accumulation. These experiments assess agreement with the chosen solver,
not independent validation of the underlying physical flow.
\end{abstract}
\keywords{fluid simulation, regression, convolutional neural networks, cross-validation}
\maketitle
\section{Research Questions}
\textbf{RQ1:} How accurately does the CNN predict the next fluid state compared
with persistence and local ridge regression? We compare final-test errors after
selecting configurations on development data.
\textbf{RQ2:} How do hidden width and learning rate affect CNN performance?
We compare configurations using identical trajectory folds.
\textbf{RQ3:} How does error accumulate when predictions are repeatedly fed back
into each method? We measure error at fixed rollout horizons from unseen starts.
@@LEARNINGRQ@@

\section{Dataset}
@@DATASET@@
Initial conditions contain random or spatially smoothed velocities; pressure
starts at zero and the moving lid is applied by the first solver step. These
synthetic velocities are not guaranteed divergence-free. We reject nonfinite
reference trajectories and screen the advective CFL number against one, in
addition to the solver's diffusion timestep check. These checks do not establish
numerical accuracy. Full seeds, physical configurations, descriptive channel
statistics, and content checksums accompany the reproducible dataset.

\section{Baselines and Experimental Setup}
Persistence predicts $S_{t+1}=S_t$. Our local ridge regressor~\cite{ridge} predicts a signed
state increment from a shared $5\times5$ stencil. The CNN uses two $3\times3$
convolutions, a hidden ReLU, and a linear signed output, also predicting an
increment. Both learned methods receive normalized $u,v,p$, spatial coordinates,
viscosity, and lid speed; both have a $5\times5$ receptive-field extent. Ridge is
implemented using NumPy normal equations and a PyTorch convolution; the CNN uses
PyTorch~\cite{pytorch} and Adam~\cite{adam}. Ridge's pixel sampling budget is reported below.

@@SETUP@@
For field $c$, the primary metric is
$E=\frac13\sum_{c\in\{u,v,p\}}\mathrm{RMSE}_c/\sigma_c$, where $\sigma_c$ is
computed from training input states only. We also retain per-field MAE and MSE.
Training minimizes mean normalized squared error; checkpoint selection minimizes
validation $E$. Raw pressure and velocity errors are reported separately because
they have different units. CV variation describes fold variability, not a
confidence interval. @@REPO@@

% Deliberately start the results on page 2. The packaging command checks the
% compiled page count; edit/condense the narrative if the draft exceeds 2 pages.
\clearpage
\section{Results and Discussion}
\subsection{RQ1: Comparative Prediction Performance}
@@CVTABLE@@
@@TESTTABLE@@
@@COMPARISON@@

\subsection{RQ2: Hyperparameter Sensitivity}
@@SENSITIVITY@@
\begin{figure}[ht]
\centering
\includegraphics[width=\columnwidth]{hyperparameters.png}
\caption{Development CV mean and fold standard deviation for CNN configurations.}
\end{figure}

\subsection{RQ3: Autoregressive Error}
@@ROLLOUT@@
\begin{figure}[ht]
\centering
\includegraphics[width=\columnwidth]{rollout.png}
\caption{Final-test rollout error. Nonfinite failures, if present, are recorded in the results tables.}
\end{figure}
@@LEARNINGRESULT@@

\subsection{Limitations and Reproducibility}
These results apply to the saved synthetic parameter distribution and numerical
solver~\cite{solver}. Testing new starts does not establish extrapolation to unseen physical
regimes, grids, or real fluids. We use one initialization seed per fold; fold
variation includes different training partitions and seeds. Reasons for model
differences require additional controlled experiments. Complete measurements,
per-trajectory errors, checkpoints, histories, splits, and environment versions
are saved alongside this report. @@TIMING@@

\begin{thebibliography}{4}
\bibitem{ridge} A. E. Hoerl and R. W. Kennard. 1970. Ridge Regression:
Biased Estimation for Nonorthogonal Problems. \textit{Technometrics} 12(1), 55--67.
\url{https://doi.org/10.1080/00401706.1970.10488634}.
\bibitem{adam} D. P. Kingma and J. Ba. 2014. Adam: A Method for Stochastic
Optimization. \url{https://arxiv.org/abs/1412.6980}.
\bibitem{pytorch} PyTorch contributors. PyTorch documentation.
\url{https://docs.pytorch.org/docs/stable/}.
\bibitem{solver} Project numerical solver, \texttt{src/simulation/solver.py}.
Numerical example attribution retained in source:
\url{https://www.youtube.com/watch?v=BQLvNLgMTQE}.
\end{thebibliography}
\end{document}
